\section*{Sammanfattning}
\addcontentsline{toc}{section}{Sammanfattning}

\begin{itemize}
  \item \code{Frame} är en \textbf{aggregate}: ingen konstruktor, ingen destruktor, inga metoder.
    Det håller aggregate-initieringen läsbar och gör klassen till ren data.
  \item \code{MaxDataLen} är \code{static constexpr} --- typ, räckvidd, synlig i debuggern, och
    samma tal som arrayens storlek.
  \item \code{isValid()} är en fri funktion, och regeln uttrycks \textbf{en gång}. Hårdvaran
    validerar inget, så drivrutinen måste --- och bara den.
  \item Du skriver inga tester i kursen. Du kör dem, läser dem och får dem att passera.
  \item Testning avgör kursen därför att hårdvaran varken finns eller är din förrän näst sist.
  \item Läs ett fallerande testfall i ordningen testnamn $\rightarrow$ Arrange $\rightarrow$
    Assert, och öppna din egen kod sist.
  \item Testsviten är oftare en exaktare kravspecifikation än prosan.
\end{itemize}
